\documentclass{ltjsbook}
\usepackage{docmute} 
\usepackage{../../myStyle} 

\begin{document}
\chapter{因子分析モデルの行列表現}
\label{x08_eigen}

\section{固有値と固有ベクトル}
今回は正方行列にみられる面白い特徴である，\keyterm{固有値}{こゆうち}{eigenvalue}と\keyterm{固有ベクトル}{こゆうべくとる}{eigenvector}についての話から始めます。
ある正方行列$\bm{A}$，列ベクトル$\bm{x}$，スカラー$\lambda$が次のような関係にあった時，$\lambda$を固有値，$\bm{x}$を固有ベクトルと言います。

\[ \bm{Ax} = \lambda\bm{x} \]

一見すると，$\bm{x}$が両辺に入っていますから，$\bm{A}$が$\lambda$に置き換わった等式に見えます。しかし一方は行列で，他方はスカラーです。こんな奇妙なことが本当にあるのでしょうか？
具体的な数値例をみてみましょう。

\[
	\bm{A} = \begin{pmatrix}
		1 & 6 \\ 2 & 5
	\end{pmatrix}
	,\bm{x} = \begin{pmatrix}
		1 \\ 1
	\end{pmatrix}
\]

を例にします。この時次の関係が成り立ちます。

\[
	\bm{Ax} =
	\begin {pmatrix}
	1 & 6 \\
	2 & 5 \\
	\end{pmatrix}
	\begin{pmatrix}
		1 \\ 1
	\end{pmatrix}
	=
	\begin{pmatrix}
		7 \\ 7
	\end{pmatrix}
	=7
	\begin{pmatrix}
		1 \\ 1
	\end{pmatrix}
	=7 \bm{x}
\]

確かに成立する組み合わせがありますね。この行列$\bm{A}$に対して，7が固有値，$(1,1)$が固有ベクトルになっています。
また，実はこの行列$\bm{A}$については，-1も固有値であり，そのときの固有ベクトルは$(-3,1)$も固有ベクトルです。

この固有値分解こそ，因子分析を元とする多変量解析の中心的な数学原理なのです。多変量解析の世界においては，分散共分散行列やそれを標準化した相関行列など，変数同士の関係を分析のスタートにおくのでした。これらの行列は正方行列ですから，その固有値や固有ベクトルを計算することで正方行列の特徴を別の視点から分解して考えられるようになります。

\subsection{固有値の特徴}
この固有値の数学的特徴は色々あるのですが，データ分析をする上で重要な点を押さえておきましょう。

固有値の特徴として，\textbf{固有値の総和が正方行列の対角要素の総和に合致する}，というのがあります。数式で表現すると，次のようになります。

\[ \sum_{i=1}^N \lambda_i = trace(\bm{A}) = \sum_{i=1}^N a_{ii} \]

ここで$a_{ij}$は行列$\bm{A}$の要素であり，$a_{ii}$は$i$行$i$列目，つまり対角要素です。
この正方行列$\bm{A}$のサイズは$N$で，対角要素の総和を特に\keyterm{トレース}{とれーす}{trace}といい$trace(\bm{A})$と表します。
それが固有値の総和とイコールになる，ということを表しています。サイズ$N$の正方行列からは固有値が$N$個算出できることがわかっており，それを全て足し合わせたものがトレースと同じになっているのですね。先ほどの例で言えば，$\bm{A}$のトレースは$1+5=6$で，固有値の総和は$7-1=6$であり，確かにこの関係が成立していることがわかります。

分散共分散行列のトレースは，分散の総和を意味します。項目同士の関係を表した行列であれば，分散はその項目から得られる情報の大きさであり，それを総和するということは，その調査研究・項目群から得られる情報の総和であると言ってもいいでしょう。
相関行列のトレースは，対角項に入っているのが$r_{ii}=1.0$ですから，項目の数と一致します。1つの項目の情報量を1.0に基準化してN項目分の情報がある，ということを表しています。

固有値都行列の関係は冒頭で示したように，$\bm{Ax}=\lambda \bm{x}$であり，正方行列の特徴をスカラーにしてしまうというものです。得られる$N$個の固有値は，元の正方行列のエッセンスをスカラーにして表現しているわけです。ここで固有値の大きさの順に並べ替えて，順に$\lambda_1,\lambda_2 \cdots \lambda_n$としたとしましょう。ここで固有ベクトルの大きさを$\bm{xx'}=1$となるように選んでやれば，元の行列を次のように書き換えることができます。

\[
	\bm{Axx'} = \lambda \bm{xx'}
\]

ここから，次のように行列$\bm{A}$を分解して行くことができます。

\[
	\bm{A}= \lambda_1 \bm{x}_1 \bm{x}_1' + \lambda_2 \bm{x}_2 \bm{x}_2' + \cdots \lambda_m \bm{x}_m \bm{x}_m'= \sum \lambda_{i} \bm{x}_{i} \bm{x}_i'
\]

これらをまとめると，固有値は元の行列の対角要素の情報を表したものであり，全体の情報量はそのままに重要度の順に並べ替えたものである，といえるでしょう。
これこそ\keytermL{因子分析}{いんしぶんせき}で取り出そうとしている因子であり，固有値の大きさはその因子の重要度である，ということがわかります！

\section{固有値と固有ベクトルを求める}
ここで少し数学の方に話を戻して，固有値と固有ベクトルの計算方法を考えましょう。
元の式を書き換えて次のような方程式を考えます。

\[ (\bm{A} - \lambda\bm{I})\bm{x} = \bm{0} \]

固有ベクトルは$\bm{x}=\bm{0}$すなわち全部ゼロであれば当然成り立ちますから(自明な解といいます)，これは除外することにします($\bm{x}\neq \bm{0}$)。
行列の表現は連立方程式の解を求めることと同じなのでした。$\bm{A} - \lambda\bm{I}$を連立方程式の係数行列だと考えれば，それが逆行列を持つと左辺にそれをかけてしまえば全部ゼロの答えになってしまいますから，そうでない答えを求めるには，この係数行列が逆行列を持たないことが重要です。

さて，この授業の中では説明してきませんでしたが，方程式が解を持つかどうかを決定する計算方法があります。これを\keyterm{行列式}{ぎょうれつしき}{determinant}といい\footnote{行列式は数値であり，解が求まるかどうか決定determinantする，という意味なのに，日本語訳はなぜか「式」といいます。変なの。}，この値がゼロでなければその方程式は解を持つ，ということがわかっています。説明しなかったのは，この値を求める計算がとても面倒だからで，詳しくは線形代数のテキストにお任せするとして\footnote{例えば\citet{村上正康2016-03-01}の第3章をみてください。}，ここでは簡便のために$2\times2$方程式の行列について紹介します。

$2 \times 2$の係数行列，$\bm{P}=\begin{pmatrix} a&b\\c&d \end{pmatrix}$の逆行列は次の式で求められることがわかっています。
\[ \bm{P}^{-1}= \frac{1}{ad-bc}\begin{pmatrix}d & -b \\ -c &a \end{pmatrix} \]
この式から考えると，$ad - bc$のところが0になるとこの計算はできませんから，逆行列が存在しないことになります。この$ad-bc$にあたるところが行列式であり，$|\bm{P}|$とか$det(\bm{P})$のように表します。
$ad-bc$が0でなければ方程式は解けるのですが，今回の場合は解けると自明になってしまうので困ります。今回は$ad-bc=0$でなければならないのです。つまり一般的に書くと次のようになります。

\[ | \bm{A} - \lambda \bm{I} | = 0 \]

$\bm{A} = \begin{pmatrix} a & b \\ c & d \end{pmatrix} $とすると，この式は次のようになります。

\[   \begin{vmatrix} a-\lambda & b \\ c & d - \lambda \end{vmatrix} = 0 \]
\[ (a-\lambda)(d-\lambda) - bc = 0 \]

この方程式を特に\keytermL{固有方程式}{こゆうほうていしき}といいますが，これを解いてやれば良いことになりますね。具体的に$\begin{pmatrix}1&6\\2&5\end{pmatrix}$の例で計算してみましょう。
\[ |\begin{pmatrix}1-\lambda & 6 \\ 2 & 5-\lambda\end{pmatrix}| =0\]
\[ (1-\lambda)(5-\lambda) - 12 = 0 \]
\[ \lambda^2 - 6\lambda -7 = 0 \]
\[ (\lambda-7)(\lambda+1)=0\]
ここから$\lambda=7,-1$が得られますね。

では固有ベクトルはどうなるでしょうか。固有値7の例で計算してみます。
\[ \begin{pmatrix} 1 & 6 \\ 2 & 5 \end{pmatrix}\begin{pmatrix} x\\y\end{pmatrix} = 7 \begin{pmatrix} x\\y\end{pmatrix} \]
\[ x + 6y = 7x \to 6x = 6y \to x=y \]
\[ 2x+5y = 7y \to 2x = xy \to x=y \]

あれっ？ なんじゃこりゃ？ と思った人もいるかもしれません。でもこれ，間違いではありません。実は固有ベクトルは大きさが定まっておらず，今回の例で言えば$x=y$つまり$x:y=1:1$の関係であればあらゆる値が成立してしまうのです。
固有ベクトルが$(1,1)$でも$(2,2)$でも$(100,100)$でも，$\bm{Ax}=\lambda \bm{x}$の関係に影響しませんから，普通はベクトルの長さ(\keyterm{ノルム}{のるむ}{norm})を1.0に規格化するという方策が取られます\footnote{ノルムとは，要素の二乗和の平方根，$\sqrt{x_1^2 + x_2^2 \cdots \x_n^2}$のことです。}。
先ほど「固有ベクトルを適当な大きさに選んでやれば」というような表現をしましたが，それはベクトルの大きさはいくらでもいいからできることなのですね。

さてここでみたように，$2\times2$の方程式であれば固有方程式を解くことはできるのですが，行列のサイズがどんどん大きくなると一般的に解けなくなって行くことは想像にかたくないと思います。実際我々は正方行列として，項目の情報が詰まった分散共分散行列とか相関行列を使いますから，それが2項目しかないなんてことはなくて，もっともっと大きなサイズになります。そうすると計算機を使って近似的に答えを求めて行くことになります。

\section{固有値と固有ベクトルの幾何学的意味}

固有値，固有ベクトルについて，今度は違う側面から見直してみましょう。

ある正方行列から固有値$\lambda$と固有ベクトル$\bm{a}$が得られたとします。このベクトル$\bm{a}$の全ての要素を定数$c$倍したベクトル$\bm{b}=c\bm{a}$を考えると，これもやはり同じ関係が成り立ちます。
\begin{equation}
	\bm{A}\bm{b}=  \lambda c \bm{a} =\lambda \bm{b}
\end{equation}
先ほどの計算でも明らかになりましたが，固有ベクトルの値は絶対的なものではなく，要素間の相対的大きさを反映しているに過ぎないのでしたね。

さてこれを幾何学的に，図形として考えてみましょう。要素が2つのベクトルは，2次元座標に表現できます。ベクトル$\bm{x}=(x,y)$という座標を表しているというわけです。固有ベクトルも要素が2つであれば，座標で表現できます。先ほどの，要素を$c$倍しても固有ベクトルとしての性質は変わらない，という話は，「固有ベクトルは大きさに意味はなく，方向を表したもの」ということになります。では何の方向を指し示しているのでしょうか。

固有値と固有ベクトルの話の最初にあった，$\bm{Ax}=\lambda\bm{x}$というのを見直してみましょう。$\bm{x}$がなんらかの座標を表していると考えると，それに正方行列をかけるとはどういう意味でしょうか。次の計算式を見てください。
\begin{equation}
	\bm{Ax}=
	\begin{pmatrix}
		2 & 0 \\
		0 & 3
	\end{pmatrix}
	\begin{pmatrix}
		1 \\
		2
	\end{pmatrix}
	=
	\begin{pmatrix}
		2 \\
		6
	\end{pmatrix}
\end{equation}

これをみると，座標$\bm{x}=\begin{pmatrix}1\\2\end{pmatrix}$に$\bm{A}$をかけたことで，座標が$\begin{pmatrix}2\\6\end{pmatrix}$に変わった，と見ることもできますね。このように，ある座標が別の座標に移ることを特に「変換」と呼びます\footnote{より一般的にいうと，以下のようになります。:集合$X$の各元$x$に集合$Y$の元$f(x)$を対応させる対応$f$のことを, 集合$X$から集合$Y$への写像(mapping)，関数(function)，あるいは変換(transformation)という。}。つまり\textbf{正方行列はなんらかの変換を施すもの}だ，と考えることができます。

今回の例では，行列$\bm{A}$には次のような性質があります。
\begin{equation}
	\begin{pmatrix}
		2 & 0 \\
		0 & 3
	\end{pmatrix}
	\begin{pmatrix}
		1 \\
		0
	\end{pmatrix}
	= 2 \begin{pmatrix} 1 \\ 0 \end{pmatrix}
\end{equation}

\begin{equation}
	\begin{pmatrix} 2 & 0\\ 0 & 3 \end{pmatrix} \begin{pmatrix} 0 \\ 1 \end{pmatrix}
	= 3 \begin{pmatrix} 0 \\ 1 \end{pmatrix}
\end{equation}

そう，お気づきのように，これは固有値・固有ベクトルです。この行列の固有値・固有ベクトルはそれぞれ$\lambda_1=2,\bm{x}_1=(1,0)$，$\lambda_2=3,\bm{x}_2=(0,1)$であることがわかります。この固有値，固有ベクトルの組み合わせをじっとみていると，面白い特徴がわかって来ます。

今回の行列から得られた2つの固有ベクトル，$(1,0)$と$(0,1)$は，2次元平面の単位ベクトルと呼ばれるものです。\textbf{2次元座標の任意の点は，これら2つのベクトルの任意の線型結合で表現できます}。座標$(a,b)$は$a\times(1,0)$と$b\times(0,1)$からなるベクトルですから，$\begin{pmatrix}1 & 0\\0 & 1\end{pmatrix}\begin{pmatrix}a\\b\end{pmatrix}=\begin{pmatrix}a\\b\end{pmatrix}$というように，です。

このように，単位ベクトルは2次元世界の基礎となる単位ともいうべきもので，ここで$(1,0)$はx座標の，$(0,1)$はy座標の基盤となるベクトルであるということができます(これを特に\textbf{基底}といいます)。つまり，正方行列$\bm{A}$から得られる固有ベクトルは，その正方行列が作る空間の基盤を明らかにするものであったのです。

では固有値はどうでしょうか？ 今回はx座標を2倍，y座標を3倍に引き延ばす変換をしたわけですが，この座標の歪み(重み)が固有値に対応していますね。つまり，固有値と固有ベクトルは新しい座標に変換する，その変換先の空間的性質を表していることになります。元の座標空間は$(1,0),(0,1)$で作られる空間ですが，変換先の空間は$(2,0),(0,3)$で作られている空間，ということになります。

つまり，正方行列は空間を変換するもの，あるいは正方行列の中に固有ベクトルを基底とした空間があるもの，ということです。

全ての正方行列に，こういった「変換」という解釈ができるのであれば，相関行列にも同様のことがいえるでしょう。相関行列は正方行列ですので，固有値分解できるのです。相関行列を固有値分解することは，相関行列の中に潜む次元(dimension)を抽出してくることです。固有ベクトル(因子負荷行列)は，正方行列によって変換される，変換先の単位ベクトルのことだったのです。そして，固有値はその次元のゆがみ(重み，重要性)という意味があったのです。

\section{因子分析の数学的理解}
さあ因子分析に戻って考えてみましょう。因子分析のモデルは次のようなものでした。
\begin{equation}
	\bm{R} = \bm{A} \bm{A} ' + \bm{D}^2
\end{equation}

ここで，左辺の正方行列を相関行列$\bm{R}$とし，$\lambda_1 \bm{xx'} = \bm{aa'}$となるようにベクトルの大きさが整えられているとすれば，次のように表現できます。
\begin{equation}
	\bm{R}=\bm{a}_1\bm{a}_1'+\bm{a}_2\bm{a}_2'+\cdots +\bm{a}_m\bm{a}_m'+\bm{d} \bm{d}'
\end{equation}

ここでは共通因子の数が$m$個だとわかっている体で分解していますが，基本的にはサイズ$N$の行列からは$N$個の固有値がずらーっと並ぶわけです。それをどこかで「共通しているのはここまで」と判断し，残りは誤差であるとしてまとめて$\bm{dd'}$にしているだけですね。このように数学的にはここからが共通因子，ここからが独自因子といった区別をすることなく，最後のひとかけらまで固有値分解を行なっているのですが，その次元の重要度でもって共通因子と誤差因子に(研究者が恣意的に)分割しているのが因子分析のやっていることなのです。

一般に，$N$よりも$m$のほうがグッと少なくなります。たとえばYG性格検査では$N=120$であり，$m$はせいぜい5から10数個です。120項目のつくる120次元空間の中で，そこに働きかけても方向の変わらない基礎的な少数の次元にのみ注目すれば，効率よく情報圧縮ができるというものです。

相関係数を固有値分解すると，その固有値は全て足し合わせるとサイズ$N$になるのでした。元のデータから計算される相関行列は，1つの項目が一単位分の情報を持っていると考えますが，固有値分解はそれを次元の重要性順に並べ替えます。固有値は項目いくつ分の重要度があるかということを表す指標だと考えることができます。どこから共通因子でどこからが誤差か，ということを考えるときに，例えば固有値が1.0よりも小さくなるようであれば，項目1つ分の情報もないのだからということで誤差因子だと判断することがあります。

ところで因子分析モデルの$\bm{A}$，因子負荷行列ですが，これは共通因子の固有ベクトルをセットで扱ったものです。つまり，$\bm{A} = \bm{a_1,a_2,\cdots,a_m}$と縦ベクトルを並べたものになっています。エレメントワイズで表現すると次のようになります。

\[
	\bm{A} = \left (
	\begin{pmatrix}  a_{11}\\a_{12}\\\vdots a_{1N} \end{pmatrix},
	\begin{pmatrix}  a_{21}\\a_{22}\\\vdots a_{2N} \end{pmatrix},
	\cdots
	\begin{pmatrix}  a_{m1}\\a_{m2}\\\vdots a_{mN} \end{pmatrix},
	\right )
\]

ここで$\bm{AA'}$の間に単位行列$\bm{I}$ を挟んでも，別に結果は変わりませんよね。単位行列はかけても変わらないのが特徴ですから，$\bm{AA'}=\bm{AIA}$です。ここでの$\bm{I}$のサイズは$m \times m$であることに注意しつつ聞いてください。

$\bm{I} = \bm{TT'}=\bm{T'T}$になるような$m$次の行列を使うと，$\bm{AA'} = \bm{ATT'A'}$の関係は保たれたままです。この$\bm{T}$はこれまた行列の座標を変えてしまう変換行列であり，これを挟むことができるということは\textbf{因子負荷量の値はなんでもあり}だということになってしまいます。この変換行列$\bm{T}$は特に\keyterm{回転行列}{かいてんぎょうれつ}{rotation matrix}とも呼ばれますが，因子分析はこのようにどんな回転でもできる，\textbf{回転に関する不定性}があるのです。あらゆる因子負荷量の組み合わせがあり得る，というのは困りものなので，なんらかの形で因子負荷行列$\bm{A}$に制約をかける必要があります。言い方を変えれば，色々な制約の中ではあっても好きな値を取ることができますから，ユーザにとって便利な基準を考えてやれば良いでしょう。因子分析では一般に，\keytermL{因子軸の回転}{いんしじくのかいてん}を行いますが，それはこうした理由からです。

ちなみに$\bm{TT'}$に
$diag(\bm{T\Phi T'})=\bm{I}_m$という
制約のある行列$\bm{\Phi}$を挟んでやっても，元の計算モデルに影響はありません。
この$\bm{\Phi}$は\keyterm{因子間相関}{いんしかんそうかん}{factor correlations}とよばれ，相関を持った因子軸の回転，すなわち\keyterm{斜交回転}{しゃこうかいてん}{oblique rotation}も色々なものが考えられています\footnote{これに対して$\bm{TT'}=\bm{I}$のような因子間相関がない(単位行列)な回転を\keyterm{直交回転}{ちょっこうかいてん}{orthogonal rotation}といいます。}。


ずいぶんややこしい話になってきました。この後は実践形式で因子分析を理解していければと思います。

\section{課題}
\paragraph{固有値問題}行列$\begin{pmatrix}8&1\\4&5\end{pmatrix}$の固有値・固有ベクトルを求めよ。


\end{document}
